\subsection{有界测度的情形}

命题 2. --- 设 \(X_{1},\ldots,X_{n},Y\) 为局部紧空间，\(\mu_{i}\) 为 \(X_{i}\) 上的有界测度（\(1\leqslant i\leqslant n\)），\(\mu\) 为各 \(\mu_{i}\) 的乘积，\(\varphi\) 为 \(\mu\)-可测映射，定义在 \(\prod_{i}X_{i}\) 上并取值于 Y。则 \(\mu_{i}\) 关于 \(\varphi\) 可卷积，并且

\[
\left\| \begin{array}{c} n \\ * \\ i = 1 \end{array} \mu_ {i} \right\| \leqslant \prod_ {i = 1} ^ {n} \| \mu_ {i} \|.
\]

若 \(\mu_{i}\) 还都是正测度，则 \(\left\| \begin{array}{c}n\\ *\\ i = 1 \end{array} \mu_i\right\| = \prod_{i = 1}^{n}\| \mu_{i}\|\)

因为 \(\mu_{i}^{\prime}=|\mu_{i}|\) 有界且 \(\|\mu_{i}^{\prime}\|=\|\mu_{i}\|\)（Ch. III, §1, No.~8, 命题 10 的推论 1）。有 \(|\mu_{1}\otimes\cdots\otimes\mu_{n}|=\mu_{1}^{\prime}\otimes\cdots\otimes\mu_{n}^{\prime}\)（Ch. III, §4, Nos. 2, 4），因此 \(\mu_{1}\otimes\cdots\otimes\mu_{n}\) 有界，并且

\[
\left\| \mu_ {1} \otimes \dots \otimes \mu_ {n} \right\| = \left\| \mu_ {1} \right\| \dots \left\| \mu_ {n} \right\|
\]

（同上，命题 4）。因此 \(\varphi\) 是 \(\mu\)-适当的（Ch. V, §6, No.~1, 注 1），也就是说，\(\mu_i\) 关于 \(\varphi\) 可卷积。有 \(\left\| \begin{array}{c} n \\ * \\ 1 \end{array} \mu_i'\right\| = \| \mu_1' \otimes \cdots \otimes \mu_n'\|\)（Ch. V, §6, No.~2, Th. 1），因而 \(\left\| \begin{array}{c} n \\ * \\ i=1 \end{array} \mu_i'\right\| = \| \mu_1'\| \cdots \| \mu_n'\|\)。最后，\(| *_{i} \mu_i| \leqslant *_{i} \mu_i'\)（No.~1, 公式 (2)），因此

\[
\left\| \ast_ {i} \mu_ {i} \right\| \leqslant \left\| \ast_ {i} \mu_ {i} ^ {\prime} \right\| = \prod_ {i = 1} ^ {n} \| \mu_ {i} \|.
\]

命题 3. --- 设 \(X_1, \ldots, X_n, Y\) 为局部紧空间，\(\varphi\) 为从 \(\prod_{i=1}^{n} X_i\) 到 \(Y\) 的连续映射。则映射

\[
(\mu , \dots , \mu_ {n}) \mapsto * _ {\varphi} (\mu_ {i})
\]

从 \(\prod_{i=1}^{n} \mathcal{M}^1(\mathrm{X}_i)\) 到 \(\mathcal{M}^1(\mathrm{Y})\) 是连续多线性映射。

这由命题 2 及 No.~1 中所述立即得到。

